% ============================================================
%  III. System Description
%  IEEE double-column format — passive / third-person voice only
%  Symbols: v_i = task node  |  u_k = UAV  |  t = decision epoch
%  All constants verified against:
%    configs/params.py · core/mobility.py · env/uav_env.py
%    core/task_shaping.py
% ============================================================

\section{System Description}
\label{sec:system}

The system under consideration consists of a swarm of $M$
heterogeneous UAVs tasked with executing a latency-sensitive
computational workload represented as a directed acyclic graph (DAG).
Unlike conventional offloading frameworks that treat the task graph
as a fixed input, the proposed approach allows the DAG topology itself
to be \emph{adaptively reshaped} at each decision epoch in response
to the dynamic resource state of the swarm.
The remainder of this section formalises the four core components of
the system: the task and shaping model, the heterogeneous swarm
representation, the mobility and communication model, and the latency
and energy consumption model.

% ----------------------------------------------------------------
\subsection{Task and Shaping Model}
\label{subsec:task_shaping}

\subsubsection*{Task DAG}

The workload is modelled as a DAG
\begin{equation}
  \mathcal{G}_T = \bigl(\mathcal{V}_T,\;\mathcal{E}_T\bigr),
  \label{eq:task_dag}
\end{equation}
where $\mathcal{V}_T = \{v_1, v_2, \ldots, v_N\}$ is the ordered set
of computational tasks and
$\mathcal{E}_T \subseteq \mathcal{V}_T \times \mathcal{V}_T$ is the
set of directed data-dependency edges.
A directed edge $(v_i, v_j) \in \mathcal{E}_T$ indicates that the
output of $v_i$ must be transmitted to the executor of $v_j$ before
$v_j$ may commence.
The acyclic structure guarantees a well-defined topological execution
order.

Each task $v_i \in \mathcal{V}_T$ is characterised by the attribute
pair
\begin{equation}
  \mathbf{x}_i^T = \bigl(c_i,\; s_i\bigr),
  \label{eq:task_attr}
\end{equation}
where $c_i > 0$ is the \emph{CPU demand} (processing units required)
and $s_i > 0$ is the \emph{local data volume} (MB).
Each edge $(v_i, v_j)$ carries a \emph{transfer volume}
$d_{ij} \geq 0$ (MB), representing the intermediate result that must
be communicated from the UAV executing $v_i$ to the UAV executing
$v_j$.
The predecessor and successor sets of $v_i$ are denoted
$\Pi(v_i) = \{v_p : (v_p, v_i) \in \mathcal{E}_T\}$ and
$\Sigma(v_i) = \{v_s : (v_i, v_s) \in \mathcal{E}_T\}$,
respectively.

\subsubsection*{Topology Shaping Operations}

The central premise of the framework is that the granularity of task
parallelism encoded in $\mathcal{G}_T$ should adapt to the
instantaneous resource state of the swarm.
Two graph-transformation primitives are defined to this end.

\smallskip
\noindent\textbf{Definition 1 (Split).}
\emph{Given} $\mathcal{G}_T = (\mathcal{V}_T, \mathcal{E}_T)$
\emph{and a target node} $v_i \in \mathcal{V}_T$
\emph{with} $|\mathcal{V}_T| < N_{\max}$,
\emph{the} \textsc{Split}$(v_i)$
\emph{operation produces}
$\mathcal{G}_T' = (\mathcal{V}_T', \mathcal{E}_T')$
\emph{as follows.}

\begin{enumerate}[label=(\roman*), leftmargin=*, itemsep=2pt]
  \item A sibling node $v_i^{+}$ is introduced.
        The attributes of $v_i$ are bisected and assigned equally to
        both the original node and its sibling:
        % Optimized for width
        \begin{equation}
          \begin{aligned}
            c_i &\leftarrow \tfrac{c_i}{2}, &\quad
            s_i &\leftarrow \tfrac{s_i}{2},\\
            c_{i^{+}} &= \tfrac{c_i}{2}, &\quad
            s_{i^{+}} &= \tfrac{s_i}{2}.
          \end{aligned}
          \label{eq:split_attr}
        \end{equation}
        A linear partition of $s_i$ is assumed for tractability;
        any decomposition synchronisation overhead is implicitly
        encapsulated within the updated CPU demand $c_i$.

  \item The updated vertex set is
        $\mathcal{V}_T' = \mathcal{V}_T \cup \{v_i^{+}\}$.

  \item The sibling inherits the complete neighbourhood of $v_i$:
        % Optimized for width
        \begin{align}
          \mathcal{E}_T' =\;& \mathcal{E}_T \nonumber\\
          &\;\cup\; \bigl\{(v_p,\, v_i^{+}) : v_p \in \Pi(v_i)\bigr\}
          \nonumber\\
          &\;\cup\; \bigl\{(v_i^{+},\, v_s) : v_s \in \Sigma(v_i)\bigr\},
          \label{eq:split_edges}
        \end{align}
        with the transfer volumes on rewired edges identical to those
        on the corresponding edges of $v_i$.
        \emph{No edge is inserted between $v_i$ and $v_i^{+}$},
        so the two sub-tasks are topologically independent and may
        execute concurrently on separate UAVs.
\end{enumerate}

The \textsc{Split} operation increases graph cardinality by one
($|\mathcal{V}_T'| = |\mathcal{V}_T| + 1$) and halves the per-sub-task
CPU demand, thereby exposing additional parallelism that is beneficial
when the swarm contains under-utilised, high-capacity UAVs.

\smallskip
\noindent\textbf{Definition 2 (Merge).}
\emph{Given two distinct nodes} $v_i, v_j \in \mathcal{V}_T$
\emph{satisfying the acyclicity guard}
\begin{equation}
  \nexists\;\text{directed path}\;v_i \leadsto v_j
  \;\;\text{and}\;\;
  \nexists\;\text{directed path}\;v_j \leadsto v_i,
  \label{eq:cycle_guard}
\end{equation}
\emph{the} \textsc{Merge}$(v_i, v_j)$
\emph{operation produces}
$\mathcal{G}_T' = (\mathcal{V}_T', \mathcal{E}_T')$
\emph{as follows.}

\begin{enumerate}[label=(\roman*), leftmargin=*, itemsep=2pt]
  \item The attributes of $v_j$ are absorbed into $v_i$:
        \begin{equation}
          c_i \leftarrow c_i + c_j, \qquad
          s_i \leftarrow s_i + s_j.
          \label{eq:merge_attr}
        \end{equation}

  \item Let $\mathcal{E}_j = \{e \in \mathcal{E}_T : v_j \in e\}$.
        Edges formerly incident to $v_j$ are redirected to $v_i$,
        without creating duplicates or self-loops:
        % Optimized for width
        \begin{align}
          \mathcal{E}_T' =\;& \bigl(\mathcal{E}_T \setminus \mathcal{E}_j\bigr)
          \nonumber\\
          &\;\cup\; \bigl\{(v_p,\, v_i) :
            v_p \in \Pi(v_j),\; v_p \neq v_i,
          \nonumber\\
          &\qquad\qquad\quad (v_p, v_i) \notin \mathcal{E}_T\bigr\}
          \nonumber\\
          &\;\cup\; \bigl\{(v_i,\, v_s) :
            v_s \in \Sigma(v_j),\; v_s \neq v_i,
          \nonumber\\
          &\qquad\qquad\quad (v_i, v_s) \notin \mathcal{E}_T\bigr\}.
          \label{eq:merge_edges}
        \end{align}

  \item Node $v_j$ is removed:
        $\mathcal{V}_T' = \mathcal{V}_T \setminus \{v_j\}$.
\end{enumerate}

The acyclicity guard \eqref{eq:cycle_guard} is both necessary and
sufficient to guarantee that $\mathcal{G}_T'$ remains a valid DAG.
The \textsc{Merge} operation decreases graph cardinality by one
($|\mathcal{V}_T'| = |\mathcal{V}_T| - 1$) and, critically, eliminates
the inter-UAV data transfer $d_{ij}$ that would otherwise be incurred
on the link connecting the executors of $v_i$ and $v_j$.
This elimination is particularly impactful when that link is
bandwidth-constrained due to large inter-UAV separation, as
characterised by the model in Section~\ref{subsec:mobility}.

% ----------------------------------------------------------------
\subsection{Heterogeneous UAV Swarm Model}
\label{subsec:uav_model}

\subsubsection*{UAV Fleet and Resource Graph}

The UAV swarm is represented as a set of $M$ heterogeneous agents
$\mathcal{U} = \{u_1, u_2, \ldots, u_M\}$.
The swarm topology is captured by a \emph{resource graph}
\begin{equation}
  \mathcal{G}_R = \bigl(\mathcal{U},\; \mathcal{L}\bigr),
  \label{eq:resource_graph}
\end{equation}
where $\mathcal{L} = \{(u_k, u_l) : k \neq l\}$ forms a fully
connected edge set, reflecting the assumption that any pair of active
UAVs may establish a communication link subject to the channel model
in Section~\ref{subsec:mobility}.

\subsubsection*{Static Hardware Profiles}

Each UAV $u_k$ is characterised by a fixed hardware profile
$\Gamma_k = (\phi_k,\, \beta_k^{(0)})$, where $\phi_k > 0$ is the
\emph{nominal CPU capacity} (processing units under full battery) and
$\beta_k^{(0)} \in (0, 1]$ is the \emph{initial battery level}
(normalised fraction of total energy capacity).
Heterogeneity is a defining characteristic of the fleet: UAVs
differ across compute capability, wireless interface bandwidth, and
initial energy reserves, reflecting the diversity of commercial
multi-rotor platforms deployed in mobile edge-computing missions.

\subsubsection*{Active UAV Set and Dynamic Resource State}

At each epoch $t$, the subset of non-depleted UAVs is defined as
\begin{equation}
  \mathcal{U}_a(t) = \bigl\{u_k \in \mathcal{U} :
    \beta_k(t) > 0 \bigr\},
  \label{eq:active_set}
\end{equation}
with cardinality $M_a(t) = |\mathcal{U}_a(t)|$.
Only members of $\mathcal{U}_a(t)$ participate in task execution and
wireless communication; depleted UAVs are excluded from scheduling
and contribute zero bandwidth to any link.

The instantaneous state of the resource graph is encoded by a
per-UAV feature matrix
$\mathbf{X}^R(t) \in [0, 1]^{M \times 3}$,
whose $k$-th row is the normalised feature vector
\begin{equation}
  \mathbf{x}_k^R(t) =
  \Bigl[\,
    \tilde{\phi}_k(t)/\phi^{\max},\;
    \bar{B}_k(t)/B_{\max},\;
    \beta_k(t)\,
  \Bigr],
  \label{eq:res_features}
\end{equation}
where $\tilde{\phi}_k(t)$ is the effective (battery-derated) CPU
capacity defined in Section~\ref{subsec:latency_energy},
$\phi^{\max} = \max_{k} \phi_k$ is the fleet-wide CPU ceiling, and
$B_{\max}$ is the peak channel capacity.
The mean link bandwidth from $u_k$ to all currently active peers is
\begin{equation}
  \bar{B}_k(t) =
  \frac{1}{M_a(t) - 1}
  \sum_{\substack{u_l \in \mathcal{U}_a(t) \\ l \neq k}} B_{kl}(t),
  \label{eq:mean_bw}
\end{equation}
where $B_{kl}(t)$ is the pairwise link bandwidth defined in
Section~\ref{subsec:mobility}, and the expression is defined to be
zero when $M_a(t) = 1$.
All three components of $\mathbf{x}_k^R(t)$ are normalised to
$[0, 1]$, yielding a scale-invariant input for the graph neural
network encoder and supporting stable policy gradient estimation.

The resource graph $\mathcal{G}_R$ is thus \emph{dynamic}: as UAVs
deplete their batteries, $M_a(t)$ contracts, reducing both the
computational capacity available to the scheduler and the bandwidth
available to inter-task communication.
This degradation motivates the adaptive topology-shaping mechanism
described in Section~\ref{subsec:task_shaping} and drives the
design of the reinforcement learning agent in
Section~\ref{sec:method}.

% ----------------------------------------------------------------
\subsection{Mobility and Communication Model}
\label{subsec:mobility}

The dynamic resource state of the swarm is fundamentally governed by
UAV mobility: as UAVs move, inter-node distances change continuously,
which in turn modulates the available link bandwidth and, therefore,
the communication latency incurred during task execution.
This non-stationarity is the primary motivation for adaptive topology
shaping.

\subsubsection*{Random Waypoint Mobility Model}

UAV positions evolve in a bounded two-dimensional arena
$\mathcal{A} = [0, L]^2$ according to the \emph{Random Waypoint}
(RWP) model~\cite{rwp_model}.
Each active UAV $u_k \in \mathcal{U}_a(t)$ moves at constant cruise
speed $v_s$ toward a target waypoint $\mathbf{w}_k(t) \in \mathcal{A}$.
At discrete step $t$ with time granularity $\Delta t$, the position
is updated as
\begin{equation}
  \mathbf{p}_k(t+1) =
  \begin{cases}
    \mathbf{w}_k(t),
      & \|\mathbf{w}_k(t) - \mathbf{p}_k(t)\|_2 \leq v_s \Delta t, \\[5pt]
    \mathbf{p}_k(t)
      + v_s \Delta t \cdot \hat{\mathbf{d}}_k(t),
      & \text{otherwise,}
  \end{cases}
  \label{eq:rwp}
\end{equation}
where the unit direction vector is
% Optimized for width
\begin{equation*}
  \hat{\mathbf{d}}_k(t) =
  \frac{\mathbf{w}_k(t) - \mathbf{p}_k(t)}
       {\|\mathbf{w}_k(t) - \mathbf{p}_k(t)\|_2}.
\end{equation*}
Upon reaching $\mathbf{w}_k(t)$, a new waypoint is sampled uniformly
from $\mathcal{A}$.
Depleted UAVs ($u_k \notin \mathcal{U}_a(t)$) remain stationary and
contribute neither computation nor communication to the system.

\subsubsection*{Distance-Dependent Bandwidth Decay}

The achievable link bandwidth between two active UAVs $u_k$ and $u_l$
at epoch $t$ is governed by a \emph{distance-dependent exponential
decay} function.
Let $r_{kl}(t) = \|\mathbf{p}_k(t) - \mathbf{p}_l(t)\|_2$ denote
the Euclidean inter-UAV distance.
The bandwidth is then
% Optimized for width
\begin{equation}
  B_{kl}(t) =
  \begin{cases}
    B_{\min} + \bigl(B_{\max} - B_{\min}\bigr)\,
      \exp\!\bigl(-r_{kl}(t)/\delta_H\bigr),
    & r_{kl}(t) \leq D_c, \\[4pt]
    B_{\min},
    & r_{kl}(t) > D_c,
  \end{cases}
  \label{eq:bandwidth}
\end{equation}
where $B_{\max}$ is the peak link capacity achieved at zero
separation, $B_{\min}$ is the floor bandwidth that persists beyond
the maximum communication radius $D_c$, and $\delta_H$ is the
\emph{characteristic decay distance} governing the rate of
degradation with range.
By convention, $B_{kk}(t) = 0$ for all $k$; and $B_{kl}(t) = 0$
whenever $u_k \notin \mathcal{U}_a(t)$ or $u_l \notin
\mathcal{U}_a(t)$.

Three operating regimes emerge from \eqref{eq:bandwidth}:
a \emph{near-field} regime ($r_{kl} \ll \delta_H$) where
$B_{kl} \approx B_{\max}$;
an \emph{intermediate} regime where bandwidth decays smoothly with
the exponential kernel; and
an \emph{out-of-range} regime ($r_{kl} > D_c$) where only the
residual floor $B_{\min}$ remains available.
Because $r_{kl}(t)$ varies continuously with swarm mobility, the
bandwidth matrix $\mathbf{B}(t) \in \mathbb{R}_{\geq 0}^{M \times M}$
is time-varying, making inter-task communication latency a
non-stationary function of the swarm configuration.
This establishes a direct coupling between UAV movement,
link quality, and the DAG transformation decisions of the
shaping agent, as formalised in the following subsection.

% ----------------------------------------------------------------
\subsection{Latency and Energy Consumption Model}
\label{subsec:latency_energy}

The time-varying bandwidth matrix $\mathbf{B}(t)$ introduced in
Section~\ref{subsec:mobility} directly determines the communication
cost of each inter-task data transfer.
Combined with the energy dynamics that progressively reduce fleet
capacity, these quantities constitute the reward-relevant signals
that drive the topology-shaping policy.

\subsubsection*{Deterministic Greedy Scheduler}

Task-to-UAV assignment is performed by a \emph{mobility-aware greedy
scheduler} that is a \emph{deterministic component of the environment},
not a learned policy.
At each epoch $t$, this scheduler processes tasks in topological
order, assigning each task $v_j$ to the active UAV
$u_l \in \mathcal{U}_a(t)$ with the greatest residual CPU budget.
The reinforcement learning agent operates at a strictly higher level
of abstraction: it observes the consequences of each scheduling round
(encoded in the feedback vector $\mathbf{f}(t)$, defined below) and
selects topology-shaping actions to influence future scheduling
outcomes, without directly controlling individual task-to-UAV
assignments.

\subsubsection*{Task Execution Latency}

The \emph{computation latency} of task $v_j$ assigned to $u_l$ is
modelled as inversely proportional to the available processing
capacity:
\begin{equation}
  L_j^{\mathrm{cmp}} = \kappa \cdot \frac{c_j}{\tilde{\phi}_l(t)},
  \label{eq:latency_cmp}
\end{equation}
where $\kappa > 0$ is a hardware-dependent latency scaling coefficient
(ms per unit of demand-to-capacity ratio) that captures processor and
operating-system overheads.

The \emph{communication latency} that must be resolved before $v_j$
may begin is the aggregate of transfer times over all incoming edges:
\begin{equation}
  L_j^{\mathrm{com}}(t)
  = \sum_{v_i \in \Pi(v_j)}
    \frac{d_{ij}}{B_{\sigma(i)\,l}(t)},
  \label{eq:latency_com}
\end{equation}
where $\sigma(i) \in [M]$ denotes the index of the UAV assigned to
predecessor $v_i$ and $B_{\sigma(i)\,l}(t)$ is the instantaneous
link bandwidth from that UAV to $u_l$ as given by
\eqref{eq:bandwidth}.
This formulation makes $L_j^{\mathrm{com}}(t)$ an explicit function
of the current swarm configuration: tasks whose executors are
separated by large distances incur high communication latency,
providing a quantitative and direct incentive for the shaping agent
to invoke \textsc{Merge} on such pairs.

The \emph{finish time} of $v_j$ under the greedy schedule is
\begin{equation}
  f_j(t) = \max_{v_i \in \Pi(v_j)} f_i(t)
            + L_j^{\mathrm{com}}(t)
            + L_j^{\mathrm{cmp}},
  \label{eq:finish_time}
\end{equation}
where $f_j(t) = 0$ for all source nodes ($\Pi(v_j) = \emptyset$).
The \emph{schedule makespan} is
\begin{equation}
  T_{\mathrm{span}}(t) = \max_{v_j \in \mathcal{V}_T}\; f_j(t),
  \label{eq:makespan}
\end{equation}
and a schedule is \emph{feasible} if and only if
$T_{\mathrm{span}}(t) \leq T_D$, where $T_D$ is the system-level
end-to-end latency deadline.

\subsubsection*{Energy Depletion Model}

Battery levels are tracked as normalised fractions
$\beta_k(t) \in [0, 1]$ and are decremented after each scheduling
round according to a linear additive cost model.
The per-round battery depletion of $u_k$ is
\begin{equation}
  \Delta\beta_k(t) =
    \alpha_c\,\rho_k(t) + \alpha_t\,\tau_k(t),
  \label{eq:energy_depletion}
\end{equation}
where $\rho_k(t) \geq 0$ (CPU units) and $\tau_k(t) \geq 0$ (MB)
denote the total computation load and data volume transmitted by
$u_k$ in round $t$, respectively.
The coefficient $\alpha_c > 0$ is the energy rate per consumed CPU
unit and $\alpha_t > 0$ is the energy rate per transmitted megabyte.
The battery state is then updated as
\begin{equation}
  \beta_k(t+1) = \max\!\bigl(\beta_k(t) - \Delta\beta_k(t),\; 0\bigr).
  \label{eq:battery_update}
\end{equation}

\subsubsection*{Battery-Aware CPU Derating}

To model the reduced computing capability of energy-constrained UAVs,
the effective CPU capacity of $u_k$ is defined as a piecewise function
of its residual battery:
\begin{equation}
  \tilde{\phi}_k(t) =
  \begin{cases}
    0,
      & \beta_k(t) = 0
        \quad\text{(depleted — offline)}, \\[4pt]
    \gamma\,\phi_k,
      & 0 < \beta_k(t) < \beta_{\mathrm{low}}
        \quad\text{(low-power mode)}, \\[4pt]
    \phi_k,
      & \beta_k(t) \geq \beta_{\mathrm{low}}
        \quad\text{(nominal operation)},
  \end{cases}
  \label{eq:cpu_derating}
\end{equation}
where $\beta_{\mathrm{low}} \in (0, 1)$ is the low-battery threshold
and $\gamma \in (0, 1)$ is the CPU derating factor applied in
low-power mode.
UAVs satisfying $\beta_k(t) = 0$ are removed from the active pool
$\mathcal{U}_a(t)$: they are excluded from both task assignment and
link formation in $\mathbf{B}(t)$.
This battery-induced shrinkage of $M_a(t)$ creates a positive
feedback loop: fewer active UAVs reduce the schedulable CPU budget,
which increases per-task execution latency \eqref{eq:latency_cmp}
and may force reschedule events—all of which are captured in the
feedback signal defined below.

\subsubsection*{Scheduling Feedback Signal}

At each epoch $t$, the outcome of the scheduling round is summarised
by the six-dimensional feedback vector
% Optimized for width
\begin{equation}
  \begin{split}
    \mathbf{f}(t) = \bigl[\,
      \lambda(t),\; \sigma(t),\; \varrho(t),\\
      \bar{u}(t),\; \hat{B}_{\min}(t),\; \hat{\beta}_{\min}(t)\,
    \bigr]^{\!\top} \in \mathbb{R}^6,
  \end{split}
  \label{eq:feedback}
\end{equation}
whose components are defined as follows.
$\lambda(t) = T_{\mathrm{span}}(t)/T_D \in \mathbb{R}_{\geq 0}$ is
the \emph{normalised makespan} (latency ratio), with $\lambda > 1$
indicating a deadline miss.
$\sigma(t) \in [0, 1]$ is the \emph{task success rate}, i.e., the
fraction of tasks in $\mathcal{V}_T$ for which a feasible assignment
was found in $\mathcal{U}_a(t)$.
$\varrho(t) \in \mathbb{Z}_{\geq 0}$ is the \emph{reschedule count},
i.e., the number of tasks that required a forced assignment to an
overloaded UAV due to insufficient residual CPU.
$\bar{u}(t) \in [0, 1]$ is the \emph{mean fleet CPU utilisation}
across active UAVs.
$\hat{B}_{\min}(t) \in [0, 1]$ is the \emph{normalised minimum link
bandwidth}, defined as
$\min_{u_k, u_l \in \mathcal{U}_a,\, k \neq l} B_{kl}(t) / B_{\max}$,
which serves as an indicator of link fragmentation induced by mobility.
$\hat{\beta}_{\min}(t) \in [0, 1]$ is the \emph{minimum battery
fraction} across the entire swarm,
$\min_{u_k \in \mathcal{U}} \beta_k(t)$, signalling impending UAV
depletion.

The vector $\mathbf{f}(t)$ is concatenated with the GNN-encoded
state of $(\mathcal{G}_T, \mathcal{G}_R)$ to form the complete
observation fed to the RL agent at each epoch.
The coupling among $\mathbf{f}(t)$, $\mathbf{B}(t)$, and the shaped
graph $\mathcal{G}_T'$ constitutes the core closed-loop signal of the
proposed framework: each topology decision modifies the workload
structure, which in turn alters the scheduling outcome and hence
$\mathbf{f}(t+1)$, enabling the agent to learn multi-step adaptive
strategies as detailed in Section~\ref{sec:method}.
